\section{Calculation methodology}

\subsection{From outages to unavailability}
Over the reporting period the forced energy unavailability follows directly from
the equivalent forced outage hours:
\[
  \mathrm{FEU}\,(\%) = \frac{\mathrm{EFOH}}{\mathrm{PH}}\times 100,
  \qquad
  \mathrm{EFOH} = \sum \mathrm{EFOD} = \sum \bigl(\mathrm{AOD}\times\mathrm{ODF}\bigr).
\]

\subsection{Liquidated damages}
For each guaranteed metric with rate $r$, measured value $m$, guarantee $g$ and
cap $C$:
\[
  \mathrm{LD} = \min\!\Bigl( r \times \underbrace{(m-g)^+}_{\text{excess}},\; C \Bigr),
  \qquad (x)^+ = \max(x,0).
\]
The total LD is the sum over the FOR and FEU components, subject to the overall
cap (3--5\,\% of contract value).

\subsection{Worked example (illustrative)}
Consider a 1\,GW scheme with a guaranteed $\mathrm{FEU}=0.5\,\%$ that measures
$0.8\,\%$ over the monitoring period, at a rate of \$300\,k per 0.1\,\% excess:
\[
  \mathrm{LD}_{\mathrm{FEU}} = \$300\text{k} \times \frac{0.8-0.5}{0.1}
  = \$300\text{k}\times 3 = \$0.9\,\text{M}.
\]
If two forced outages occur against a guarantee of one, at \$0.5\,M per outage,
$\mathrm{LD}_{\mathrm{FOR}} = \$0.5\,\text{M}$. The combined \$1.4\,M is then
capped at the contractual maximum (e.g.\ 4\,\% of a \$200\,M supply value $=$
\$8\,M, so the cap does not bind here).

\begin{center}
\renewcommand{\arraystretch}{1.3}
\begin{tabular}{@{}l r r@{}}
  \toprule
  \textbf{Component} & \textbf{Basis} & \textbf{LD} \\
  \midrule
  FEU (0.3\,\% excess) & 3 $\times$ \$300\,k & \$0.90\,M \\
  FOR (1 extra outage) & 1 $\times$ \$0.5\,M & \$0.50\,M \\
  \midrule
  \textbf{Total (pre-cap)} & & \textbf{\$1.40\,M} \\
  \bottomrule
\end{tabular}
\end{center}

{\color{atGray}\footnotesize Figures are illustrative; will be replaced with project
rates, guarantees and the actual OEM supply value.}
